\chapter{The framework and its relations to existing mathematics}\label{ch:synthesis}

The preceding chapters built the programme bottom-up. This chapter states the framework as a whole and places each layer against the mathematics it resembles. Two things are kept apart throughout: what is \emph{classical} (cited, not claimed), and what is \emph{specific to this programme} (a statement about generative histories of a prefix-defined or symbolic grammar). Nothing here claims to replace or extend Shannon's theory; the relation to it is a relation of \emph{complement}, made precise in \cref{sec:shannon}. A claim-by-claim comparison, with the exact statement, the earlier mathematics, what is inherited, and the unresolved novelty questions, is kept in the file \texttt{PRIOR\_ART.md} of the package; the bibliographic status of every reference is recorded in \cref{app:sources}.

\section{The framework in one statement}\label{sec:framework}
Fix a finite set of $k$ distinguishable events and a set $A\subseteq S_k$ of admissible orderings (histories). The programme studies four objects attached to $A$ and the relations between them.
\begin{enumerate}[label=(\roman*)]
\item \emph{Multiplicity.} The counts $|A|$, $|F_t(m_0)|$ of histories, and the Shannon and Hartley quantities derived from them (\cref{ch:ancestry}).
\item \emph{Local legality.} The swap graph $G_A$ of $A$ with the colour of each swap, and the radius-$r$ balls in it (\cref{ch:swaps,ch:access}).
\item \emph{Topology.} The history complex $K_A\subseteq\partial P_k$, its Betti numbers over $\F$, and the cycles and obstructions they record (\cref{ch:topology,ch:atlas,ch:dimension}).
\item \emph{Forbidden structure.} The complement complex $Y_A$, its nerve $N[D_A]$ on the minimal forbidden prefixes, and the core of that nerve (\cref{ch:duality,ch:minimal,ch:locality}).
\end{enumerate}
The relations established are of three kinds. (a)~\emph{Determination}: $G_A$ with its colours determines $K_A$ and every $\beta_j$ (\cref{thm:swap-reconstruction}), and $K_A$ determines $Y_A$ up to the duality of \cref{thm:AD}. (b)~\emph{Separation}: histories can have equal multiplicity data and different topology (\cref{ch:separation}, witnesses W1--W6), and equal local data to a given radius and different topology (\cref{thm:radius2} and the radius-two failures). (c)~\emph{Arithmetic control}: for the SELECT-II grammar the family $\{K_{A_C(m)}\}_m$ is a divisibility filtration governed by finitely many thresholds (\cref{thm:face-threshold}), and the survival of a given hole is governed by a finite-lattice law (\cref{thm:AT1}).

The central thesis of the programme is then the following, stated as a summary and not as a theorem. \emph{Information about generative history is organised by two independent structures. One is the multiplicity of histories, to which Shannon and Hartley measures apply directly. The other is the combination of arithmetic or symbolic constraints, the topology of the admissible region, and the structure of the forbidden histories; this determines which questions about the history an observation of a given resolution can answer.} The evidence for the second structure being independent of the first is \cref{prop:witnesses,obs:equal-N}; the evidence for it being decisive for recovery is \cref{thm:obs-ladder,thm:no-radius1}, \cref{obs:preregistered,obs:hard-positives}, \cref{prop:counts-fail}. All of these are scoped to the corpus stated in the respective chapters.

\section{Shannon and Hartley}\label{sec:shannon}
Hartley's logarithmic measure of the number of possible sequences \cite{Hartley28} and Shannon's entropy \cite{Shannon48} are used in \cref{ch:ancestry,ch:separation} exactly as defined there; the chain rule, Gibbs' inequality and the data-processing facts used are the textbook ones \cite{CoverThomas}. Two points of contact are worth stating precisely.

\emph{What is inherited.} The ambiguity $A(t)=\log_2|F_t|$ of \cref{ch:ancestry} is the Hartley measure of the set of complete histories, and the conditional entropy $\Ent(M|T)$ is Shannon's. Their relation (\cref{thm:hierarchy}) is an instance of the classical inequality between Hartley and Shannon entropy of a distribution on a finite set.

\emph{What is not claimed.} We do not define a new entropy, and we do not claim that any topological quantity is an entropy. The point of \cref{ch:separation} is a \emph{negative} statement about the sufficiency of a one-number summary: on the stated comparable windows, histories-count (equivalently Hartley entropy of the uniform distribution on $A$) does not determine the Betti vector. This is a statement that the \emph{distribution} on $A$ is an incomplete description of the \emph{support} $A$, and says nothing against Shannon theory, whose subject is the distribution. The mutual-information reading of \cref{ch:predictive} (plug-in versus out-of-sample information) is likewise standard estimation practice applied to a finite corpus.

\section{Algebraic topology: duality, nerves, strong homotopy}\label{sec:algtop}
\paragraph{Alexander duality.} \Cref{thm:AD} is the classical Alexander duality for a compact, locally contractible subset of a sphere \cite{Hatcher}, applied to $|Y_A|\subset|\partial P_k|\cong S^{k-2}$ after identifying $|K_A|$ with a deformation retract of the complement. The identification (Step~1 of its proof) is a standard barycentric-subdivision argument \cite{Munkres}. The only programme-specific content is the dictionary: history complexes are complements of forbidden-history complexes in the sphere of orderings, so every statement about holes in the admissible region is a statement about cohomology of the forbidden region. Combinatorial Alexander duality for simplicial complexes on a finite vertex set \cite{BjornerTancer} is a different statement (it dualises a complex on $n$ vertices within the full simplex on the same vertex set) and is not what is used here; we cite it only as the nearest combinatorial relative. Likewise the Eagon--Reiner theorem relating the Alexander dual of a squarefree monomial ideal to linear resolutions \cite{EagonReiner} concerns ideals, not histories, and is mentioned only to prevent a conflation of the two dualities.

\paragraph{Nerve theorems.} The dead-set nerve theorem (\cref{thm:nerve}) is an instance of the classical nerve theorem for a finite cover by subcomplexes whose non-empty intersections are acyclic (or contractible) \cite{Borsuk,Bjorner}. The programme-specific step is the star-intersection \cref{lem:star}: intersections of the generating sub-complexes $Z_p$ of $Y_A$ are non-empty exactly when the corresponding forbidden prefixes are compatible, and then are cones, hence contractible. The closure of the family under intersection that makes the cover acyclic is Quillen's fibre lemma in the form used in \cref{lem:quillen} \cite{Quillen}. The statement that $Y_A$ is homotopy equivalent to a complex on the \emph{minimal forbidden prefixes} is a consequence; it is analogous to, but should not be confused with, the use of nerves of covers by maximal faces in persistent topology \cite{Carlsson09}.

\paragraph{Strong homotopy and cores.} The reduction of \cref{ch:minimal} uses strong collapses and cores of finite simplicial complexes \cite{BarmakMinian}: removal of a dominated vertex is a strong homotopy equivalence, and every finite complex has a core, unique up to isomorphism (\cref{thm:barmak}). We use these as black boxes. What is specific is the empirical \emph{size distribution} of cores of nerves $N[D_A]$ on the corpus (\cref{tab:D21}), and the corollary that, for those complexes, the core is a complete invariant for the homology of $K_A$ while no count-level statistic is (\cref{prop:counts-fail}).

\section{Lcm lattices and cellular resolutions}\label{sec:lcm}
This section records a reading of \cref{thm:face-threshold} in the language of combinatorial commutative algebra \cite{MillerSturmfels}. It is offered as an interpretation resting on a cited theorem, not as a new result, and none of the later chapters depends on it.

\paragraph{Setting.} For a prefix-defined set the vertices $w\in S_k$ of the permutohedron $P_k$ carry integer labels $K(w)$, and \cref{thm:face-threshold}(a) says that the face $F$ carries the label $K_F=\lcm_{w\in V(F)}K(w)$ and that $F\in X_C(m)$ iff $K_F\mid m$. Identify a positive integer $n$ with the monomial $x^{n}=\prod_q x_q^{v_q(n)}$ in variables indexed by primes; divisibility of integers is then divisibility of monomials, $\lcm$ is the monomial $\lcm$, and the faces of $P_k$ carry exactly the labelling convention of a labelled cell complex in the sense of Bayer and Sturmfels \cite{BayerSturmfels}. For a monomial degree $b$ the subcomplex of faces whose label divides $b$ is the complex $X_{\le b}$ of that paper; in our notation $X_{\le b}=X_C(m)$ when $b$ is the exponent vector of $m$.

\paragraph{The interpretation.} By \cite[Prop.~1.1]{BayerSturmfels}, the cellular complex of a labelled complex supports a free resolution of the monomial module generated by its vertex labels if and only if every $X_{\le b}$ is acyclic over the base field (with the conventions of that paper for the empty complex). Applied to the labelled permutohedron $P_k$ with labels $K(w)$ this gives:
\begin{quote}
\emph{the labelled permutohedron supports a cellular resolution of $\langle x^{K(w)}:w\in S_k\rangle$ over a field $\Bbbk$ iff, for every $m$, the history complex $X_C(m)$ is $\Bbbk$-acyclic.}
\end{quote}
The converse reading is the one that matters here: a hole of $X_C(m)$, such as the anchor class of \cref{def:anchor} with $\beta_2>0$ at $m=288$, is an \emph{obstruction} to the permutohedron supporting a resolution of that module. The thresholds of \cref{ch:topology} are then the multidegrees at which the obstruction appears or disappears. We have not computed the module, its minimal resolution, or its Betti numbers, and make no claim about them. The classical lcm-lattice of the module \cite{GasharovPeevaWelker} is the finite lattice of $\lcm$'s of subsets of the generators; the finitely many distinct systems $A_C(m)$ of \cref{prop:finite-lattice}, indexed by capped exponent vectors, are of the same lcm type. We have not examined the precise relation between the two lattices. Whether the survival law \cref{thm:AT1} can be phrased as a statement about the lcm-lattice, and whether the nonvanishing of $\beta_2(X_C(m))$ gives information about the minimal free resolution of the module, are open (\cref{prob:resolution}).

\section{Symbolic dynamics}\label{sec:symbolic-dyn}
The grammars of \cref{ch:locality} are the finite-memory constraints of symbolic dynamics: a pairwise grammar is a one-step shift of finite type on the alphabet of events, restricted to words that are permutations (each letter exactly once) \cite{LindMarcus}. Two differences determine what can be transferred. First, the words are of fixed finite length and are permutations, so none of the ergodic or entropy-rate notions of symbolic dynamics (topological entropy of the shift, zeta function) apply; the object is the \emph{finite} set of admissible words and its topology. Second, the constraint-locality result (\cref{thm:pairwise}) is an assertion about $H_2$ of the permutation complex, which has no counterpart in the usual theory. The sharpness example (\cref{prop:triple}) shows that moving from a pairwise (memory-one, 2-letter forbidden words) to a triple constraint (memory two, 3-letter forbidden words) creates $\beta_2>0$ at $k=6$. In shift-space language, forbidden words of length $2$ cannot create such a hole on six letters and forbidden words of length $3$ can. We stress again that the theorem is proved for $k=6$ only.

\section{Persistent homology}\label{sec:persistence}
The family $m\mapsto X_C(m)$ is a filtration of simplicial (indeed polytopal) complexes indexed by the divisibility lattice, equivalently by the exponent vectors $(v_q(m))_q\in\N^r$. A one-parameter filtration yields a persistence module whose barcode is an isometry-stable invariant \cite{ELZ02,Carlsson09}; a filtration by $\N^r$ with $r\ge 2$ is a multi-parameter persistence module, for which no complete discrete barcode exists in general \cite{CarlssonZomorodian}. The analogy to the present setting is useful but limited.
\begin{itemize}
\item \emph{What corresponds.} Births are the thresholds of \cref{thm:face-threshold}; the exact change in Betti numbers across an inclusion is the rank--nullity identity \cref{thm:face-threshold}(c), which is the elementary identity behind every persistence computation. Survival of a fixed class (\cref{thm:AT1}) is the statement that a fixed class has a certain support in the parameter lattice.
\item \emph{What does not.} We do not compute barcodes or any multi-parameter invariant, we do not use stability, and the atlas of \cref{ch:atlas} classifies cycles by their survival \emph{laws} (a finite lattice-valued invariant of the cycle) rather than by persistence intervals. The relation to the rank invariant of multi-parameter persistence is not established.
\end{itemize}

\section{Combinatorial reconstruction}\label{sec:reconstruction}
The permutohedron is a simple polytope. A classical theorem of Blind and Mani-Levitska, with a short proof by Kalai, states that a simple polytope is determined by its graph, in the sense that the graph determines the face lattice \cite{BlindManiLevitska,Kalai88}. Ziegler's book is the standard reference for the polytope facts used \cite{Ziegler}, and Bj\"orner and Brenti for the Coxeter-group facts (parabolic cosets, the weak order) \cite{BjornerBrenti}.

\Cref{thm:swap-reconstruction} is a reconstruction statement of a different kind. The classical theorem applies to the \emph{entire} graph of a simple polytope. Here the graph is an arbitrary \emph{induced subgraph} of the graph of $P_k$, for which no reconstruction can hold without extra data (a bare induced subgraph forgets which swaps are of which kind); with the swap colours, the faces of $K_A$ are recovered by a component count that uses only the cardinalities of parabolic subgroups. We do not claim that the proof technique is new. The colours are not decorative: \cref{prop:uncoloured-fails} exhibits, for the anchor content, two sources whose uncoloured swap graphs are isomorphic with equal Betti numbers and $f$-vectors, yet which differ in whether the anchor class survives; only the labelled arithmetic data separate them. That is a statement about the anchored state and not about the Betti numbers themselves, which \cref{thm:swap-reconstruction} shows are functions of the coloured graph; whether uncoloured induced subgraphs of $P_k$ can differ in $\beta_2$ is not addressed here.

\section{What is new, and what is not}\label{sec:novelty}
\Cref{tab:novelty} lists the main results by the kind of contribution. The word ``instance'' means that the statement is a specialisation of a cited theorem; ``dictionary'' means that the proof is a translation into the language of the programme; ``finding'' means a corpus-scoped computational observation with exact verification; ``construction'' means a definition and proved properties not taken from the literature that we located.
\begin{table}[ht]\centering\footnotesize
\setlength{\tabcolsep}{4pt}
\begin{tabular}{@{}p{3.6cm}p{2.3cm}p{6.1cm}@{}}
\toprule
Result & Kind & Closest earlier mathematics\\\midrule
\cref{thm:swap-reconstruction} & dictionary & face lattice of simple polytopes from the graph \cite{BlindManiLevitska,Kalai88}\\
\cref{thm:AD} & instance & Alexander duality \cite{Hatcher}\\
\cref{thm:nerve} & instance & nerve theorem, Quillen closure \cite{Borsuk,Bjorner,Quillen}\\
\cref{thm:barmak}, \cref{cor:core-suffices} & instance & strong homotopy \cite{BarmakMinian}\\
\cref{thm:face-threshold} & dictionary & lcm labelling \cite{BayerSturmfels,MillerSturmfels}\\
\cref{thm:AT1}, atlas & construction, finding & none located; compare persistence \cite{Carlsson09}\\
\cref{ch:dimension} (lift calculus, $J_X$) & construction & none located\\
\cref{thm:radius2} (radius-two identity) & construction & elementary counting in bipartite graphs\\
\cref{thm:pairwise} & finding (computer-assisted, $k=6$) & compare shifts of finite type \cite{LindMarcus}\\
\cref{ch:separation,ch:predictive} & findings & Shannon, Hartley \cite{Shannon48,Hartley28,CoverThomas}\\
\bottomrule
\end{tabular}
\caption{Kind of contribution. ``None located'' means that the authors searched and did not find an earlier statement, not that none exists. The unresolved novelty questions are listed in \texttt{PRIOR\_ART.md}.}\label{tab:novelty}
\end{table}
